% ============================================================
% ch33.tex —— 尾声：数学之网
% ============================================================
\chapter*{尾声\quad 数学之网}
\addcontentsline{toc}{chapter}{尾声\quad 数学之网}

三百多页走完，最后做一件事：\textbf{收网}。回望这张网，你会发现六个篇章从来不是六条平行线，而是一张反复打结的网——每一次打结处，都是两个世界的互相照亮。把结一个个指给你看，是这一章全部的内容。

\section*{六条网线}

\textbf{网线一：余数的接力（篇一 $\to$ 篇四 $\to$ 篇六）。}《孙子算经》的"物不知数"（第 5 章）$\to$ 费马小定理与 RSA（第 6 章）$\to$ 欧拉的齿轮（第 22 章）$\to$ $\zeta$ 函数与黎曼猜想（第 24 章）$\to$ 椭圆曲线密码（第 32 章）。一条两千年的接力：同一类"余数问题"，从军营点兵跑到国际结算。\textbf{同余是全书跑得最远的一件工具}——它出发时只是除法竖式的零头。

\textbf{网线二：抽象的反哺（篇一 $\leftrightarrow$ 篇二）。}第 6 章用"幂迟早回头"硬啃出费马小定理的直觉版，立下欠条；第 9 章拉格朗日的盘子三行还清欠款。\textbf{抽象不是逃离具体，是具体问题的期货市场}——投入一次证明，处处收获利息。同一幕重演了三次：第 5 章"$Z_p$ 里人人有倒数"的实验，第 10 章由裴蜀等式转正；第 5 章的"魔法数"，第 11 章在多项式世界长成拉格朗日插值。

\textbf{网线三：刘徽的刀（篇三 $\to$ 篇四）。}割圆术（第 14 章）$\to$ 极限与比赛规则 $\to$ 导数（第 15 章）$\to$ 积分（第 17 章）$\to$ $\mathrm{e}$ 与 $\ln$（第 19 章）$\to$ $\mathrm{li}(x)$ 素数流量图（第 23 章）。公元 263 年的一把刀，一路切到了 1896 年素数定理的证明现场；而它最后的一次挥刀（第 32 章椭圆周长积分）恰好回到自己的起点——\textbf{切圆的刀，最后切的是"圆为什么特殊"本身}。

\textbf{网线四：蚂蚁的胜利（篇五）。}曲率（第 25 章）$\to$ 测地线（第 26 章）$\to$ 内在弯曲（第 27 章）$\to$ 宇宙的形状（第 28 章）。一条主线索贯穿全篇：\textbf{不离开表面，也能知道世界的弯曲}——这也正是人类认知宇宙的唯一处境。蚂蚁在第 27 章量出球面的那天，宇宙学的方法论就奠基了。

\textbf{网线五：账本精神（贯穿全书）。}判别式管相切（第 29 章）、切点两票（第 30 章）、循环点的公共住址（第 31 章）、弦切加法的三票（第 32 章）——\textbf{几何问题最终都是记账问题}，而账本的语言是代数。贝祖定理一次结清全部历史欠账；"次数"从记号升格为权力。反过来，账本也要看几何的脸色：第 31 章的齐次坐标之所以发明，正是因为纯代数的账（平行线欠一票）逼着几何开出新大陆。

\textbf{网线六：大合龙（第 32 章）。}椭圆曲线上长出群（篇二 $\cap$ 篇六）；它的密码学回到篇一的数论（离散对数）；它的 BSD 猜想回到篇四的 $\zeta$ 家族（$L$ 函数在 $s=1$ 的零点阶 $=$ 秩）；它的名字与椭圆周长回到篇三的积分。一条曲线，网收四角——全书六篇，在最后一条曲线上打了总结。

\section*{三张空白}

地图上还有空白，这不丢人——\textbf{知道空白在哪，才算真的有地图}：

\begin{itemize}
  \item \textbf{孪生素数猜想}（第 4 章）：相差 $2$ 的素数对是否有无穷多？张益唐 2013 年的突破把间隔压到 $246$，最后 $244$ 的路程还在走；
  \item \textbf{黎曼猜想}（第 24 章）：$\zeta$ 的零点全在临界线上？十万亿个数值证据、百分之四十的理论覆盖、零个完整证明——悬赏一百万美元；
  \item \textbf{BSD 猜想}（第 32 章）：椭圆曲线的秩能否由 $L$ 函数读出？同样在千禧年悬赏名单上。
\end{itemize}

注意这三张空白的分布：一张在素数（最古老的整数问题），一张在分析（最抽象的连续机器），一张在代数几何与数论的交界。它们不是三个孤岛——朗兰兹纲领（二十世纪数学最大胆的构想）主张：\textbf{数论、代数几何、分析这三片大陆，共享一套隐藏的"字典"}，任何一方的定理都能翻译成另两方的语言。上面三张空白，全都压在这套字典的要害上。下一个结网的人，可能是你。

\section*{给你的一封回信}

你在书里"发明"过的每个土名字，大学课本里都有正式名字：落脚点叫极限，瞬间变化率叫导数，堆积叫积分，体检四项叫群公理，盘子叫陪集，弯度叫曲率，画家平面叫射影平面，第三点传送门叫椭圆曲线的群律。\textbf{你缺的从来不是天赋，只是一张术语对照表}——它就在附录 A 里，两页纸，两小时能对完。对完暗号你再翻任何一本大学教材，会发现每一页都只是在讲你做过的事，只是穿了正装。

你更需要知道的是：你在本书里反复练习的那个动作——\textbf{遇到不懂的东西，给它起个名字，规定它怎么动，然后检查自洽}——就是从欧几里得到黎曼的两千三百年里，所有数学家的全部动作。没有别的动作了。欧几里得造了余数和素数的户口（第 1--4 章），《孙子算经》造了同余方程组的解法（第 5 章），伽罗瓦造了群（第 7--9 章），刘徽造了极限（第 14 章），牛顿莱布尼茨造了微积分（第 15--18 章），黎曼造了流形与 $\zeta$ 的复地图（第 24、28 章）——\textbf{他们和你用的是同一张发明许可证}，就是序章发你的那张。

烧掉的书已经重造完毕。这一次，每一块砖你都知道从哪来、为什么这么砌。\textbf{下一本，轮到你写。}
